Vision-language models (VLMs) have become the default tool for training-free visual anomaly detection (VAD), yet single-pass prompting fails silently on domains where the VLM's world-knowledge priors conflict with the task-specific anomaly definition. We study how to add structure on top of a VLM for \emph{cross-domain} VAD and present a training-free main system together with an optional semi-supervised controller extension.

\textbf{Main finding.} On \textbf{CrossDomainVAD-12} ($12$ domains, $1{,}418$ test items, spanning industrial, medical, remote-sensing, infrastructure, 3D, and road-scene sources), we introduce \textbf{AnomalyClaw}: an autonomous refutation agent whose single per-item invocation runs, \emph{in parallel on the same VLM backbone}, (i) a descriptor-free Direct call and (ii) a multi-turn tool-using trajectory, and returns their unweighted average. Across three VLMs, the ensemble is positive-in-mean on all three and statistically significant at $95\%$ on two: GPT-5.4 macro AUROC $0.731\!\to\!0.772$ ($+4.01$~pp, stratified paired bootstrap $95\%$ CI $[+2.58,+5.50]$, $P(\Delta\!>\!0){=}1.000$); SeedVL $0.678\!\to\!0.738$ ($+5.99$~pp, $[+4.32,+7.66]$, $P{=}1.000$); Qwen3.5-VL-27B $0.714\!\to\!0.732$ ($+1.69$~pp, $[-0.08,+3.51]$, $P{=}0.971$) --- positive in mean but the $95\%$ CI just touches zero on this third backbone.

\textbf{Failure-mode robustness.} The agent is \emph{not} individually better than Direct on every backbone. On Qwen3.5-VL-27B the agent alone is significantly \emph{worse} than Direct ($-4.46$~pp, $95\%$ CI $[-7.61,-1.55]$, $P{=}0.001$), yet the always-on parallel-Direct branch still produces a positive ensemble gain. This failure-mode robustness is the load-bearing property: an agent recipe that ensembles against its own Direct inside the same call can be dropped onto a new VLM backbone without first verifying that the agent alone is stronger than Direct. The agent and Direct make complementary errors on \emph{every} backbone we tested, even when the agent is weak.

\textbf{Semi-supervised controller extension.} On Qwen3.5-VL-27B --- the backbone where the ensemble's gain is smallest --- we study whether a learned arbitration layer can close the gap. \textbf{The controller extension} adds a Controller VLM that reads both branches' score--rationale pairs and a short rulebook retrieved per $(\mathrm{domain}, \mathrm{category})$. Rules are built offline from $40$ dev labels per domain ($480$ labels total, no inference-time supervision). The controller with its full rulebook reaches $0.7539$, a $+2.05$ pp gain over the passive ensemble baseline of $0.7333$ (paired bootstrap $95\%$ CI $[+0.41,+3.67]$, $\hat p{=}0.995$). A branch-frozen ablation (no-rules, shuffled-domain rules, meta-only, full) shows that the controller mechanism alone does not help ($-1.06$ pp), wrong-domain rules are indistinguishable from no rules ($-0.30$ pp over no-rules, $\hat p{=}0.315$), and correct-domain rule content is what carries the gain (full vs shuffled $+3.39$ pp, CI $[+1.87,+4.79]$, $\hat p{=}1.000$). The controller is a single-backbone case study; transfer to GPT-5.4 and SeedVL is left to future work.

\textbf{Takeaway.} For training-free cross-domain VAD, the right architectural primitive is not a pure autonomous agent but an \emph{ensemble-aware} agent whose Direct-branch is always invoked. When a small dev-label budget is available, a controller-level rulebook extension adds a further significant improvement on backbones where the ensemble gain is marginal. We release AnomalyClaw (refutation agent + parallel Direct) and its optional controller extension, the CrossDomainVAD-12 benchmark ($1{,}418$ items across 12 domains with calibration/dev/test splits), and full paired-bootstrap analysis artifacts for every headline comparison.
